% GENERATED FILE. Edit canonical lessons, then run npm run generate:books.
% canonical-source-hash: fnv1a64:501e93d08993eb23
% canonical-lessons: HI-C510-randa, HI-C510-cheni, HI-C510-ghisna, HI-C510-thokna, HI-C510-chedna

\chapter{In the Carpenter's Shop}
\label{ch:hi-in-the-carpenter-s-shop}
\hlchaptermodality{510}

\begin{chapteropening}
\textbf{By the end of this chapter:} \emph{I can describe how wood is shaped and finished.}

Complete the last lesson of chapter 510: I can describe how wood is shaped and finished.
\end{chapteropening}

\section[\texorpdfstring{\hi{रंदा} (randā)}{रंदा (randā)}]{\hi{रंदा} (randā) — a carpenter's plane}

\label{lesson:HI-C510-randa}



\begin{quote}
Before the new one: say the Hindi for high-handedness, then the Hindi for a slogan.
\end{quote}

\subsection*{You'll want to know: \hi{रंदा}}
\textbf{\hi{रंदा}} — \emph{randā} — \textquotedblleft{}a carpenter's plane\textquotedblright{}.

\begin{quote}
\hi{बढ़ई} \hi{ने} \hi{रंदे} \hi{से} \hi{दरवाज़ा} \hi{चिकना} \hi{किया।}
\end{quote}

\begin{quote}
\emph{Baṛhaī ne rande se darvāzā cikanā kiyā.}
\end{quote}

\begin{quote}
The carpenter made the door smooth with a plane.
\end{quote}

Masculine.

\begin{cousinweb}
\textbf{\hi{रंदा}} (\emph{randā}) is usually traced to Persian \emph{randa}, \textquotedblleft{}a plane (the tool)\textquotedblright{}.
\end{cousinweb}

\begin{grammarlens}[title={What you've built}]
One more word for describing what you lived and felt.
\end{grammarlens}

\subsection*{Your turn}
\begin{itemize}
\raggedright
  \item \emph{Say it:} \emph{randā}
  \item \emph{Say it:} \emph{randā}, once more
  \item \emph{Say it:} \emph{Baṛhaī ne rande se darvāzā cikanā kiyā.} — The carpenter made the door smooth with a plane.
  \item \emph{Recall:} say the Hindi for high-handedness, then the Hindi for a slogan, then say \emph{randā} again
\end{itemize}

\begin{tcolorbox}[breakable,colback=teal!4,colframe=teal!35!black,title={Before you move on}]
What does \hi{रंदा} mean? (\textquotedblleft{}A carpenter's plane\textquotedblright{}.) Say it once more.
\end{tcolorbox}
\section[\texorpdfstring{\hi{छेनी} (chenī)}{छेनी (chenī)}]{\hi{छेनी} (chenī) — a chisel}

\label{lesson:HI-C510-cheni}



\begin{quote}
Before the new one: say the Hindi for a slogan, then the Hindi for a carpenter's plane.
\end{quote}

\subsection*{You'll want to know: \hi{छेनी}}
\textbf{\hi{छेनी}} — \emph{chenī} — \textquotedblleft{}a chisel\textquotedblright{}.

\begin{quote}
\hi{छेनी} \hi{और} \hi{हथौड़े} \hi{से} \hi{वह} \hi{लकड़ी} \hi{में} \hi{फूल} \hi{बना} \hi{रहा} \hi{है।}
\end{quote}

\begin{quote}
\emph{Chenī aur hathauṛe se vah lakṛī meṁ phūl banā rahā hai.}
\end{quote}

\begin{quote}
With a chisel and a hammer he is carving flowers into the wood.
\end{quote}

Feminine.

\begin{cousinweb}
\textbf{\hi{छेनी}} (\emph{chenī}) is a tadbhava, inherited from Sanskrit \emph{chēdana}, \textquotedblleft{}cutting\textquotedblright{} (CDIAL 5066).
\end{cousinweb}

\begin{grammarlens}[title={What you've built}]
One more word for describing what you lived and felt.
\end{grammarlens}

\subsection*{Your turn}
\begin{itemize}
\raggedright
  \item \emph{Say it:} \emph{chenī}
  \item \emph{Say it:} \emph{chenī}, once more
  \item \emph{Say it:} \emph{Chenī aur hathauṛe se vah lakṛī meṁ phūl banā rahā hai.} — With a chisel and a hammer he is carving flowers into the wood.
  \item \emph{Recall:} say the Hindi for a slogan, then the Hindi for a carpenter's plane, then say \emph{chenī} again
\end{itemize}

\begin{tcolorbox}[breakable,colback=teal!4,colframe=teal!35!black,title={Before you move on}]
What does \hi{छेनी} mean? (\textquotedblleft{}A chisel\textquotedblright{}.) Say it once more.
\end{tcolorbox}
\section[\texorpdfstring{\hi{घिसना} (ghisnā)}{घिसना (ghisnā)}]{\hi{घिसना} (ghisnā) — to rub, to wear down}

\label{lesson:HI-C510-ghisna}



\begin{quote}
Before the new one: say the Hindi for a carpenter's plane, then the Hindi for a chisel.
\end{quote}

\subsection*{You'll want to know: \hi{घिसना}}
\textbf{\hi{घिसना}} — \emph{ghisnā} — \textquotedblleft{}to rub, to wear down\textquotedblright{}.

\begin{quote}
\hi{बढ़ई} \hi{ने} \hi{लकड़ी} \hi{को} \hi{पत्थर} \hi{पर} \hi{घिसा।}
\end{quote}

\begin{quote}
\emph{Baṛhaī ne lakṛī ko patthar par ghisā.}
\end{quote}

\begin{quote}
The carpenter rubbed the wood on a stone.
\end{quote}

Past, with \hi{ने} (he or she …): \hi{उसने} \hi{छुरी} \hi{पत्थर} \hi{पर} \hi{घिसी}, \emph{usne churī patthar par ghisī}.

\begin{cousinweb}
\textbf{\hi{घिसना}} (\emph{ghisnā}) is a tadbhava, inherited from Sanskrit \emph{gharṣati}, \textquotedblleft{}rubs\textquotedblright{} (CDIAL 4450).
\end{cousinweb}

\begin{grammarlens}[title={What you've built}]
One more word for describing what you lived and felt.
\end{grammarlens}

\subsection*{Your turn}
\begin{itemize}
\raggedright
  \item \emph{Say it:} \emph{ghisnā}
  \item \emph{Say it:} \emph{ghisnā}, once more
  \item \emph{Say it:} \emph{Baṛhaī ne lakṛī ko patthar par ghisā.} — The carpenter rubbed the wood on a stone.
  \item \emph{Recall:} say the Hindi for a carpenter's plane, then the Hindi for a chisel, then say \emph{ghisnā} again
\end{itemize}

\begin{tcolorbox}[breakable,colback=teal!4,colframe=teal!35!black,title={Before you move on}]
What does \hi{घिसना} mean? (\textquotedblleft{}To rub, to wear down\textquotedblright{}.) Say it once more.
\end{tcolorbox}
\section[\texorpdfstring{\hi{ठोकना} (ṭhoknā)}{ठोकना (ṭhoknā)}]{\hi{ठोकना} (ṭhoknā) — to knock in; to drive (a nail)}

\label{lesson:HI-C510-thokna}



\begin{quote}
Before the new one: say the Hindi for a chisel, then the Hindi for to rub, to wear down.
\end{quote}

\subsection*{You'll want to know: \hi{ठोकना}}
\textbf{\hi{ठोकना}} — \emph{ṭhoknā} — \textquotedblleft{}to knock in; to drive (a nail)\textquotedblright{}.

\begin{quote}
\hi{उसने} \hi{दीवार} \hi{में} \hi{एक} \hi{कील} \hi{ठोकी।}
\end{quote}

\begin{quote}
\emph{Usne dīvār meṁ ek kīl ṭhokī.}
\end{quote}

\begin{quote}
He drove a nail into the wall.
\end{quote}

Past, with \hi{ने} (he or she …): \hi{उसने} \hi{मेज़} \hi{में} \hi{कील} \hi{ठोकी}, \emph{usne mez meṁ kīl ṭhokī}.

\begin{cousinweb}
\textbf{\hi{ठोकना}} (\emph{ṭhoknā}) is a tadbhava, inherited from \emph{*ṭhōkk}, \textquotedblleft{}knock\textquotedblright{}, a form CDIAL reconstructs from its descendants (CDIAL 5513).
\end{cousinweb}

\begin{grammarlens}[title={What you've built}]
One more word for describing what you lived and felt.
\end{grammarlens}

\subsection*{Your turn}
\begin{itemize}
\raggedright
  \item \emph{Say it:} \emph{ṭhoknā}
  \item \emph{Say it:} \emph{ṭhoknā}, once more
  \item \emph{Say it:} \emph{Usne dīvār meṁ ek kīl ṭhokī.} — He drove a nail into the wall.
  \item \emph{Recall:} say the Hindi for a chisel, then the Hindi for to rub, to wear down, then say \emph{ṭhoknā} again
\end{itemize}

\begin{tcolorbox}[breakable,colback=teal!4,colframe=teal!35!black,title={Before you move on}]
What does \hi{ठोकना} mean? (\textquotedblleft{}To knock in; to drive (a nail)\textquotedblright{}.) Say it once more.
\end{tcolorbox}
\section[\texorpdfstring{\hi{छेदना} (chednā)}{छेदना (chednā)}]{\hi{छेदना} (chednā) — to bore; to pierce}

\label{lesson:HI-C510-chedna}



\begin{quote}
Before the new one: say the Hindi for to rub, to wear down, then the Hindi for to knock in.
\end{quote}

\subsection*{You'll want to know: \hi{छेदना}}
\textbf{\hi{छेदना}} — \emph{chednā} — \textquotedblleft{}to bore; to pierce\textquotedblright{}.

\begin{quote}
\hi{बढ़ई} \hi{ने} \hi{कील} \hi{से} \hi{लकड़ी} \hi{छेदी।}
\end{quote}

\begin{quote}
\emph{Baṛhaī ne kīl se lakṛī chedī.}
\end{quote}

\begin{quote}
The carpenter bored through the wood with a nail.
\end{quote}

Past, with \hi{ने} (he or she …): \hi{उसने} \hi{काग़ज़} \hi{छेदा}, \emph{usne kāġaz chedā}.

\begin{cousinweb}
\textbf{\hi{छेदना}} (\emph{chednā}) is a tadbhava, inherited from Sanskrit \emph{chidra}, \textquotedblleft{}a hole\textquotedblright{} (CDIAL 5043). CDIAL gives the Hindi verb as \textquotedblleft{}to bore\textquotedblright{}.
\end{cousinweb}

\begin{grammarlens}[title={What you've built}]
One more word for describing what you lived and felt.
\end{grammarlens}

\subsection*{Your turn}
\begin{itemize}
\raggedright
  \item \emph{Say it:} \emph{chednā}
  \item \emph{Say it:} \emph{chednā}, once more
  \item \emph{Say it:} \emph{Baṛhaī ne kīl se lakṛī chedī.} — The carpenter bored through the wood with a nail.
  \item \emph{Recall:} say the Hindi for to rub, to wear down, then the Hindi for to knock in, then say \emph{chednā} again
\end{itemize}

\begin{tcolorbox}[breakable,colback=teal!4,colframe=teal!35!black,title={Before you move on}]
What does \hi{छेदना} mean? (\textquotedblleft{}To bore; to pierce\textquotedblright{}.) Say it once more.
\end{tcolorbox}
